% Generated draft: requires mathematical and source review.
\documentclass[11pt]{article}
\usepackage[margin=25mm]{geometry}
\usepackage{fontspec}
\setmainfont{Latin Modern Roman}
\usepackage{amsmath,amssymb,mathtools,unicode-math}
\setmathfont{Latin Modern Math}
\usepackage{enumitem,needspace,xurl,hyperref}
\setlength{\parindent}{0pt}
\setlength{\parskip}{5pt}
\setlength{\emergencystretch}{3em}
\newcommand{\N}{\mathbb{N}}
\newcommand{\Z}{\mathbb{Z}}
\newcommand{\Q}{\mathbb{Q}}
\newcommand{\R}{\mathbb{R}}
\newcommand{\C}{\mathbb{C}}
\newcommand{\F}{\mathbb{F}}
\newcommand{\A}{\mathbb{A}}
\newcommand{\PP}{\mathbb{P}}
\newcommand{\E}{\mathbb{E}}
\newcommand{\eps}{\varepsilon}
\newcommand{\dd}{\,\mathrm{d}}
\begin{document}
\section*{UnsolvedMath ID: 701208}
\textbf{Source problem code:} catalogue entry
\subsection{Definitions and mathematical statement}
Let \(G\) be a group with identity \(1\). Use the commutator convention \([x,y]=x^{-1}y^{-1}xy\), and define \([x,{}_1 g]=[x,g]\) and \([x,{}_{n+1}g]=[[x,{}_n g],g]\) for integers \(n\geq 1\). Under the usual left Engel convention, define
\[
E(G)=\{g\in G:\ \forall x\in G\ \exists n=n(x,g)\geq 1\text{ such that }[x,{}_n g]=1\}.
\]
The integer \(n\) may depend on both \(x\) and \(g\). In particular, no bound uniform in \(x\) is required for an element to belong to \(E(G)\); bounded Engel elements are included as well. The supplied record does not explicitly specify left versus right Engel elements. The formulation below uses the displayed convention, subject to verification of the source's intended convention.

A group has no elements of order \(2\) if \(x^2=1\) implies \(x=1\). It is torsion-free if, for every \(x\in G\) and every integer \(m\geq 1\), the equality \(x^m=1\) implies \(x=1\). A group is right-orderable if it admits a strict total order \(<\) such that
\[
x<y\ \Longrightarrow\ xz<yz\qquad(x,y,z\in G).
\]
A group is bi-orderable if it admits a strict total order invariant under multiplication on both sides, equivalently an order satisfying
\[
x<y\ \Longrightarrow\ axb<ayb\qquad(a,b,x,y\in G).
\]
The source uses the phrase ``linearly ordered groups'' without defining its compatibility with multiplication. In part (d), this is interpreted as bi-orderable groups, as is customary in this context; the intended meaning requires source verification.

\textbf{Questions (V. V. Bludov, 2006; Kourovka Notebook 16.16).} Is each of the following universally quantified assertions true?
\begin{enumerate}
\item[(a)] For every group \(G\) with no elements of order \(2\), the set \(E(G)\) is a subgroup of \(G\).
\item[(b)] For every torsion-free group \(G\), the set \(E(G)\) is a subgroup of \(G\).
\item[(c)] For every right-orderable group \(G\), the set \(E(G)\) is a subgroup of \(G\).
\item[(d)] For every bi-orderable group \(G\), the set \(E(G)\) is a subgroup of \(G\), subject to the qualification about ``linearly ordered'' above.
\end{enumerate}
The groups in these questions are not assumed finite or finitely generated.

\textbf{Supplied background qualification.} The record states that examples exist of \(2\)-groups in which the product of two Engel elements is not an Engel element. Here a \(2\)-group means a group in which every element has finite order a power of \(2\); no finiteness of the group is imposed. The source describes the two elements as ``(unboundedly) Engel,'' in the context of allowing Engel conditions without a uniform bound. It does not identify these examples or explicitly clarify whether each factor must fail every bounded Engel condition. This background assertion is preserved as supplied and is not independently verified here.
\subsection{Short English statement}
Bludov asks whether the Engel elements form a subgroup in each of four classes: groups without involutions, torsion-free groups, right-orderable groups, and linearly ordered groups. The draft uses left Engel elements and interprets the last class as bi-orderable groups; both conventions require confirmation from the source. The Engel length may depend on the element being tested. The supplied registry labels the problem open, and its background reports counterexamples among \(2\)-groups.
\subsection{Sources}
\begin{itemize}
\item Local UnsolvedMath registry, permanent ID \(701208\), exact source code \texttt{catalogue entry}, title ``Kourovka Notebook Problem 16.16.'' The supplied record attributes the problem to V. V. Bludov, with proposed year \(2006\), and labels it \texttt{open}. Its creation and update timestamps are both \texttt{2026-10-16T14:00:00Z}; these are registry timestamps, not independent mathematical verification.
\item \emph{The Kourovka Notebook (Unsolved Problems in Group Theory)}, problem 16.16, originating in the 16th issue (2006), as identified in the supplied background with the 21st issue edition (2026), arXiv identifier \texttt{1401.0300v48}. The collection description names E. I. Khukhro and V. D. Mazurov as editors. This bibliographic attribution is supplied by the record; the cited edition has not been consulted for this draft.
\item Supplied upstream record ID \(701208\), revision \texttt{d371ee917ac9647fbb06cd6fc529deec6f3dda82}. No separate references for the reported \(2\)-group examples, literature-source links, or literature-check date are supplied.
\end{itemize}
\end{document}
